% TOP_PROBLEM: {"id":30007526,"problem_number":"LOCAL-30007526","title":"Measuring the neutral interest rate","category_id":21,"category":{"id":21,"name":"economics","display_name":"Economics","description":"Open economic theory statements, published research agendas, and proposed scoped specifications; question provenance and historical priority are qualified per record.","slug":"economics","order_index":21,"created_at":"2026-10-08T00:00:00Z"},"status":"provenance_pending","external_url":null,"legacy_ids":[],"published":false,"statement":"Determine whether current neutral and long-run natural real interest rates can be separately identified from demand shocks and potential-output revisions.","economics":{"original_id":15,"field":"Inflation and monetary policy","kind":"identification","formalization_basis":"proposed_specification","source_status":"not_verified","priority_status":"not_established","priority_note":"No explicit primary open-problem statement matching this catalogue item was verified; supporting research is not evidence of first formulation.","earliest_verified_reference":null,"current_reference":{"authors":["Maurice Obstfeld"],"title":"Natural and Neutral Real Interest Rates: Past and Future","year":2023,"url":"https://www.nber.org/papers/w31949","doi":"10.3386/w31949","locator":"Abstract, primary NBER indexed record","evidence_summary":"The abstract distinguishes a long-run natural rate from a short-run inflation-neutral rate and warns that one may not reliably guide the other. This is conceptual support; an explicit open identification statement was not verified in the full text."},"formalization_note":"New identification specification. The supporting source distinguishes the concepts but an exact primary open statement was not verified."},"statusQualification":"Provenance pending: an explicit published statement of this exact open question was not verified. This is a catalogue-authored candidate specification, not a certified open conjecture. New identification specification. The supporting source distinguishes the concepts but an exact primary open statement was not verified.","statusReviewedAt":"2026-10-08"}
% FIRST_POSTED: 2026-10-08
% ENTRY_KIND: statement_only
% !TeX program = lualatex
\documentclass[11pt]{article}
\usepackage{fontspec}
\usepackage[a4paper,margin=25mm]{geometry}
\usepackage{amsmath,amssymb,mathtools}
\usepackage{xurl}
\usepackage[hidelinks]{hyperref}
\newcommand{\sref}[2]{\hyperlink{source:#1:#2}{[S#2]}}
\setlength{\emergencystretch}{3em}
\begin{document}
% problemId:30007526
\section*{Measuring the neutral interest rate}\label{30007526}
\subsection{Definitions and mathematical statement}
In a specified small macroeconomic state-space model, let \(x_t\) be the output gap, \(r_t\) the ex ante real policy rate, and \(r_t^*\) the latent contemporaneous neutral rate. Assume
\[
x_t=E_tx_{t+1}-\sigma^{-1}(r_t-r_t^*)+d_t,\qquad
\pi_t=\beta E_t\pi_{t+1}+\kappa x_t+u_t,
\]
where \(\sigma,\kappa>0\), \(\beta\in(0,1)\), \(d_t\) is an additional demand disturbance, and \(u_t\) a cost disturbance. Specify observed proxies, measurement errors, transition laws, and survey information. Define the long-run natural rate \(\bar r\) as the steady-state flexible-price real rate under stated demographic, productivity, and fiscal fundamentals; it need not equal \(r_t^*\).
Determine whether \(r_t^*\), \(\bar r\), and their difference are identified from the proposed information set. In particular, exhibit or eliminate observational equivalence between shifts in \(r_t^*\), demand disturbances, and revisions to potential output. Develop confidence sets accounting for weak identification and model uncertainty, then assess their real-time stability across data vintages and policy regimes. A smooth filtered latent series is insufficient evidence of identification. The target concerns reliable distinctions between current inflation-neutral conditions and long-run equilibrium in this declared model family, rather than measurement of an observable universal interest-rate constant.

\par\textbf{Formulation and scope.} New identification specification. The supporting source distinguishes the concepts but an exact primary open statement was not verified.

\par\textbf{Open-question provenance.} An explicit published open-question statement was not verified. Sources below are supporting literature and must not be treated as a first listing.

\par\textbf{Historical priority.} No explicit primary open-problem statement matching this catalogue item was verified; supporting research is not evidence of first formulation.

\subsection{Short English statement}
Determine whether current neutral and long-run natural real interest rates can be separately identified from demand shocks and potential-output revisions.

\subsection{Sources}
\begin{itemize}\raggedright
\item[\textnormal{[S1]}]\hypertarget{source:30007526:1}{} Maurice Obstfeld. 2023. Natural and Neutral Real Interest Rates: Past and Future. Abstract, primary NBER indexed record.
\url{https://www.nber.org/papers/w31949}.
DOI: 10.3386/w31949.
{\small\textit{Supporting or current literature; see evidence qualification. The abstract distinguishes a long-run natural rate from a short-run inflation-neutral rate and warns that one may not reliably guide the other. This is conceptual support; an explicit open identification statement was not verified in the full text.}}
\end{itemize}
\end{document}
